\section{FMEA DE PRODUTO E DE PROCESSO (ENTREGA 4)}\label{sec:fmea}

% Arquivo gerado por scripts/gerar_fmea.py. Não edite à mão: altere o script e rode de novo.

O FMEA entra no projeto detalhado como parte da avaliação dos sistemas, subsistemas e componentes. Para cada função do produto e cada operação do processo, levantam-se os modos de falha, os efeitos e as causas, pontuam-se a severidade (S), a ocorrência (O) e a detecção (D), e o produto dos três índices, o número de prioridade de risco (NPR), ordena as ações \cite{rozenfeld2006gestao}. A equipe usou escalas de 1 a 10 (Tabela~\ref{tab:fmea_escalas}) e tratou todo modo de falha com NPR de 100 ou mais e todo modo com severidade de 9 ou mais, qualquer que fosse o NPR.

\begin{table}[htbp]
    \centering
    \caption{Escalas de severidade, ocorrência e detecção}
    \label{tab:fmea_escalas}
    \footnotesize
    \renewcommand{\arraystretch}{1.3}
    \begin{tabularx}{\textwidth}{@{}
        c
        >{\RaggedRight\arraybackslash}X
        >{\RaggedRight\arraybackslash}p{3.2cm}
        >{\RaggedRight\arraybackslash}X@{}}
        \toprule
        \textbf{Índice} & \textbf{Severidade} & \textbf{Ocorrência} & \textbf{Detecção} \\
        \midrule
        1 a 2 & Ninguém percebe & Remota & Detectada em todas as unidades \\
        3 a 4 & Incômodo, sem perder a compra & Rara & Teste de todas as unidades \\
        5 a 6 & A compra termina no caixa, ou há custo para a loja & Ocasional & Amostragem ou inspeção visual \\
        7 a 8 & A peça paga não sai, ou a não paga é levada & Frequente & Controle indireto \\
        9 a 10 & Furto sem registro, cobrança errada ou risco às pessoas & Quase certa & Sem controle, aparece no uso \\
        \bottomrule
    \end{tabularx}
    \source{Elaborado pelos autores (2026).}
\end{table}

As Tabelas~\ref{tab:fmea_produto} e~\ref{tab:fmea_processo} trazem as duas análises. O FMEA de produto percorre as funções dos sistemas da estrutura revisada, e o de processo, as operações das fichas e do plano macro que podem gerar aquelas falhas.

\begin{landscape}
\begingroup
\scriptsize
\setlength{\tabcolsep}{3pt}
\renewcommand{\arraystretch}{1.2}
\begin{longtable}{@{}
    >{\RaggedRight\arraybackslash}p{0.6cm}
    >{\RaggedRight\arraybackslash}p{2.0cm}
    >{\RaggedRight\arraybackslash}p{2.6cm}
    >{\RaggedRight\arraybackslash}p{2.4cm}
    c
    >{\RaggedRight\arraybackslash}p{3.0cm}
    c
    >{\RaggedRight\arraybackslash}p{2.6cm}
    c c
    >{\RaggedRight\arraybackslash}p{4.6cm}
    c c c c@{}}
    \caption{FMEA de produto do Tag\&Go}\label{tab:fmea_produto}\\
    \toprule
    & & & & & & & & & & & \multicolumn{4}{c}{\textbf{Após a ação}} \\
    \cmidrule(l){12-15}
    \textbf{Nº} & \textbf{Função} & \textbf{Modo de falha} & \textbf{Efeito} & \textbf{S} & \textbf{Causa} & \textbf{O} & \textbf{Controle atual} & \textbf{D} & \textbf{NPR} & \textbf{Ação} & \textbf{S} & \textbf{O} & \textbf{D} & \textbf{NPR} \\
    \midrule
    \endfirsthead
    \toprule
    & & & & & & & & & & & \multicolumn{4}{c}{\textbf{Após a ação}} \\
    \cmidrule(l){12-15}
    \textbf{Nº} & \textbf{Função} & \textbf{Modo de falha} & \textbf{Efeito} & \textbf{S} & \textbf{Causa} & \textbf{O} & \textbf{Controle atual} & \textbf{D} & \textbf{NPR} & \textbf{Ação} & \textbf{S} & \textbf{O} & \textbf{D} & \textbf{NPR} \\
    \midrule
    \endhead
    \bottomrule
    \endfoot
    D01 & Reter a peça & O pino sai com menos de 120~N & Peça levada sem pagamento & 8 & Embreagem fora da especificação ou pino revenido na fundição & 4 & Arrancamento por amostragem no Posto~4 & 5 & \textbf{160} & Dureza e extração do pino na inspeção da garra (Tabela~\ref{tab:proc_garra}) & 8 & 2 & 3 & \textbf{48} \\
    \addlinespace
    D02 & Reter a peça & A garra sai dos mancais da tampa & Peça levada sem pagamento & 8 & Mancais de PC/ABS cedem a uma alavanca (M7) & 5 & Nenhum & 6 & \textbf{240} & Ensaio de alavanca na homologação do produto. Eixo metálico TG-212 como reserva & 8 & 2 & 4 & \textbf{64} \\
    \addlinespace
    D03 & Liberar a peça & A embreagem não abre após o pagamento & Cliente pagou e vai ao caixa & 7 & Folga acima do limite, ímã fraco ou prateleira empenada & 5 & Gaussímetro no recebimento. Teste de todos os subconjuntos & 3 & \textbf{105} & Gabarito FD-01, batente da prensa e plano de controle do suporte & 7 & 2 & 2 & \textbf{28} \\
    \addlinespace
    D04 & Liberar a peça & A embreagem abre sem ordem & Peça levada sem pagamento & 9 & Ímã externo forte, como o de um destacador comercial & 6 & Antifurto da saída da loja & 6 & \textbf{324} & Chave da garra informa toda abertura sem ordem, e o servidor alerta a equipe da loja & 9 & 5 & 3 & \textbf{135} \\
    \addlinespace
    D05 & Liberar a peça & O ímã não volta ao repouso & Dispositivo destravado na próxima peça & 8 & Microatuador falha ou célula com tensão baixa & 3 & Teste funcional na montagem & 5 & \textbf{120} & \textit{Firmware} só gira com tensão acima do mínimo e confirma o retorno antes de informar travado & 8 & 2 & 3 & \textbf{48} \\
    \addlinespace
    D06 & Identificar o dispositivo & O celular não lê o código 2D & Cliente vai ao caixa & 5 & Gravação gasta pelo uso & 4 & Leitura de todos os códigos no Posto~0 & 3 & \textbf{60} & Releitura e regravação no recondicionamento & 5 & 3 & 2 & \textbf{30} \\
    \addlinespace
    D07 & Comunicar com o servidor & A ordem não chega ao dispositivo & Compra termina pela via humana & 6 & Cobertura de rádio insuficiente no salão & 5 & Nenhum & 4 & \textbf{120} & Mapear a cobertura na instalação de cada loja. Tempo-limite na página de compra (E08) & 6 & 3 & 3 & \textbf{54} \\
    \addlinespace
    D08 & Suprir energia & A célula acaba antes de 168 dias (E13) & Dispositivo não libera & 7 & Corrente de espera acima da prevista & 4 & Nenhum & 5 & \textbf{140} & Corrente de espera medida na homologação. Aviso de carga baixa ao servidor & 7 & 2 & 3 & \textbf{42} \\
    \addlinespace
    D09 & Suprir energia & A célula aquece ou é danificada & Risco de fogo & 10 & Indução aquece peças metálicas, ou esmagamento & 2 & Proteção da célula & 5 & \textbf{100} & Célula IEC 62133-2 e sensor de temperatura que corta a recarga & 10 & 1 & 3 & \textbf{30} \\
    \addlinespace
    D10 & Preservar a roupa & A garra marca o tecido & Peça devolvida, custo para a loja & 6 & Dentes agressivos ou mola forte & 4 & Nenhum & 5 & \textbf{120} & Homologar em tecido leve, malha e jeans da Riachuelo & 6 & 3 & 3 & \textbf{54} \\
    \addlinespace
    D11 & Proteger o mecanismo & A tampa é aberta por terceiros & Liberação sem pagamento & 8 & Parafusos comuns & 3 & Chave da garra & 6 & \textbf{144} & Parafusos hexalobulares e registro de abertura sem ordem & 8 & 2 & 5 & \textbf{80} \\
\end{longtable}
\endgroup
\begin{center}
    \footnotesize Fonte: Elaborado pelos autores (2026). S: severidade. O: ocorrência. D: detecção. NPR: número de prioridade de risco.
\end{center}
\end{landscape}

\begin{landscape}
\begingroup
\scriptsize
\setlength{\tabcolsep}{3pt}
\renewcommand{\arraystretch}{1.2}
\begin{longtable}{@{}
    >{\RaggedRight\arraybackslash}p{0.6cm}
    >{\RaggedRight\arraybackslash}p{2.0cm}
    >{\RaggedRight\arraybackslash}p{2.6cm}
    >{\RaggedRight\arraybackslash}p{2.4cm}
    c
    >{\RaggedRight\arraybackslash}p{3.0cm}
    c
    >{\RaggedRight\arraybackslash}p{2.6cm}
    c c
    >{\RaggedRight\arraybackslash}p{4.6cm}
    c c c c@{}}
    \caption{FMEA de processo do Tag\&Go}\label{tab:fmea_processo}\\
    \toprule
    & & & & & & & & & & & \multicolumn{4}{c}{\textbf{Após a ação}} \\
    \cmidrule(l){12-15}
    \textbf{Nº} & \textbf{Operação} & \textbf{Modo de falha} & \textbf{Efeito} & \textbf{S} & \textbf{Causa} & \textbf{O} & \textbf{Controle atual} & \textbf{D} & \textbf{NPR} & \textbf{Ação} & \textbf{S} & \textbf{O} & \textbf{D} & \textbf{NPR} \\
    \midrule
    \endfirsthead
    \toprule
    & & & & & & & & & & & \multicolumn{4}{c}{\textbf{Após a ação}} \\
    \cmidrule(l){12-15}
    \textbf{Nº} & \textbf{Operação} & \textbf{Modo de falha} & \textbf{Efeito} & \textbf{S} & \textbf{Causa} & \textbf{O} & \textbf{Controle atual} & \textbf{D} & \textbf{NPR} & \textbf{Ação} & \textbf{S} & \textbf{O} & \textbf{D} & \textbf{NPR} \\
    \midrule
    \endhead
    \bottomrule
    \endfoot
    P01 & Injeção do suporte (op.~20) & Prateleira empenada ou fora da espessura & Folga fora da tolerância & 7 & Contração do poliacetal, resfriamento desigual & 5 & Inspeção da op.~50 & 4 & \textbf{140} & Cpk mínimo de 1,33 no \textit{try-out}. Medição após 24~h & 7 & 3 & 2 & \textbf{42} \\
    \addlinespace
    P02 & Fundição da garra (op.~10 e 20) & Pino revenido ou inclinado & Retenção cai ou o pino não entra na tampa & 8 & Calor do alumínio, bucha do molde gasta & 4 & Visual & 6 & \textbf{192} & Dureza, perpendicularidade e extração na op.~70. Troca programada da bucha & 8 & 2 & 3 & \textbf{48} \\
    \addlinespace
    P03 & Colagem do ímã (op.~0.4 a 0.6) & O ímã descola em uso & Dispositivo não libera & 7 & Alojamento sujo, dose errada ou cura incompleta & 4 & Nenhum & 6 & \textbf{168} & Desengraxe, rugosidade no DE-05, cisalhamento por lote e registro da estufa & 7 & 2 & 3 & \textbf{42} \\
    \addlinespace
    P04 & Colagem do ímã (op.~0.5) & Altura do conjunto fora da tolerância & Folga fora da tolerância & 7 & Ímã fora do fundo do ninho & 4 & Teste do subconjunto & 3 & \textbf{84} & Pastilha de aço no FD-01 e medição de um braço por placa & 7 & 2 & 3 & \textbf{42} \\
    \addlinespace
    P05 & Prensagem do braço (op.~1.4) & Braço fora da altura ou solto & Folga fora da tolerância & 7 & Batente desregulado ou interferência baixa (M3) & 4 & Teste do subconjunto & 3 & \textbf{84} & Aferir o batente por turno. Torque de arraste por amostragem & 7 & 2 & 3 & \textbf{42} \\
    \addlinespace
    P06 & Montagem da célula (op.~3.3) & Célula perfurada & Risco de fogo & 10 & Ferramenta pontiaguda ou célula fora de posição & 2 & Visual & 4 & \textbf{80} & Nervuras de posição na base. Posto sem ferramenta metálica & 10 & 1 & 4 & \textbf{40} \\
    \addlinespace
    P07 & Aperto dos parafusos (op.~3.8) & Rosca espanada na torre & Tampa abre com pouco esforço & 6 & Torque acima do limite & 4 & Parafusadeira com torque & 4 & \textbf{96} & Torque definido na homologação. Aferição por turno & 6 & 2 & 3 & \textbf{36} \\
    \addlinespace
    P08 & Cadastro (op.~4.2) & Código vinculado a outra placa & Cliente paga e outro dispositivo abre & 9 & Dois dispositivos na bancada & 3 & Teste funcional & 3 & \textbf{81} & O cadastro acende o LED do dispositivo vinculado. Um por vez na bancada & 9 & 2 & 2 & \textbf{36} \\
\end{longtable}
\endgroup
\begin{center}
    \footnotesize Fonte: Elaborado pelos autores (2026). S: severidade. O: ocorrência. D: detecção. NPR: número de prioridade de risco.
\end{center}
\end{landscape}

\subsection{Pontos críticos e providências}
\label{subsec:fmea_criticos}

Dos 19 modos de falha, 13 tinham NPR de 100 ou mais, e outros 2 entraram pela severidade. A Tabela~\ref{tab:fmea_providencias} resume as providências que já constam dos documentos desta entrega e as que ficam para a homologação do produto. Com elas, a soma dos NPR cai de 2558 para 941, e só um modo continua acima de 100.

\begin{table}[htbp]
    \centering
    \caption{Providências do FMEA}
    \label{tab:fmea_providencias}
    \footnotesize
    \renewcommand{\arraystretch}{1.3}
    \begin{tabularx}{\textwidth}{@{}
        >{\RaggedRight\arraybackslash}p{2.4cm}
        >{\RaggedRight\arraybackslash}X
        >{\RaggedRight\arraybackslash}p{4.6cm}@{}}
        \toprule
        \textbf{Modos} & \textbf{Providência} & \textbf{Onde está} \\
        \midrule
        D03, P04, P05 & Gabarito que fixa a altura do ímã e batente da prensa & Seção~\ref{sec:projeto_dispositivo} e DE-07 \\
        D03, P01 & Características críticas e plano de controle do suporte & DE-03 e Seção~\ref{sec:controle_qualidade} \\
        D01, P02 & Dureza, perpendicularidade e extração do pino & Ficha da garra, op.~70 \\
        D09, P06 & Célula certificada, sensor de temperatura e nervuras de posição & Tabela~\ref{tab:ext_lista} e DE-01 \\
        D11 & Parafusos hexalobulares & Tabela~\ref{tab:ext_lista} \\
        P03 & Rugosidade do alojamento e cura em estufa & DE-05 e Tabela~\ref{tab:disp_uso} \\
        \midrule
        D02, D08, D10, P07 & Ensaios de alavanca, corrente de espera, tecido e torque & Homologação do produto \\
        D04, D05, P08 & Regras do \textit{firmware} e do \textit{software} de cadastro & Especificação da placa \\
        D07 & Mapa de cobertura de rádio e tempo-limite na página de compra & Instalação de cada loja \\
        \bottomrule
    \end{tabularx}
    \source{Elaborado pelos autores (2026).}
\end{table}

O modo que sobra é a abertura da embreagem por um ímã externo (D04). Destacadores capazes de abrir etiquetas rígidas são vendidos livremente \cite{cablematic2026xy021}, e o Tag\&Go herda essa fragilidade porque retém a peça pelo mesmo princípio. A diferença é que ele percebe a abertura: a chave da garra avisa o servidor, que distingue a abertura com ordem da abertura sem ordem. A providência melhora a detecção, mas não a ocorrência. Uma blindagem de aço sobre a embreagem reduziria a ocorrência, só que ela também enfraquece o campo do próprio ímã do dispositivo, e a equipe preferiu não adotá-la sem medir esse efeito.

Três dos riscos maiores (D01, D02 e P02) vêm das mudanças M2 e M7 do DFMA, que integraram o pino à garra fundida e trocaram o eixo metálico por mancais de plástico. As duas reduziram peças e tempo de montagem, e as duas mexem na retenção, que é a função que a loja mais cobra. Por isso o eixo metálico da versão anterior ao DFMA (TG-212) ficou como reserva, caso o ensaio de alavanca da homologação reprove os mancais.
